% ═══════════════════════════════════════════════════════════════════════════
% QR-Code für Lernpfad: Repaso general DS-03
% ═══════════════════════════════════════════════════════════════════════════

\documentclass[a4paper,11pt]{article}

\usepackage[utf8]{inputenc}
\usepackage[T1]{fontenc}
\usepackage{helvet}
\renewcommand{\familydefault}{\sfdefault}

\usepackage[margin=20mm]{geometry}
\usepackage{tikz}
\usepackage{xcolor}
\usepackage{qrcode}
\usepackage[hidelinks]{hyperref}

% Fachfarbe Spanisch
\definecolor{spanisch}{HTML}{c41e3a}

\pagestyle{empty}

% KONFIGURATION
\newcommand{\lptitel}{Repaso general (DS-03)}
\newcommand{\lpurl}{https://devbox42.github.io/sciencesim/spanisch/kl09/unidad-02-mi-nueva-casa/DS-03/LP-03-repaso-general.html}
\newcommand{\lpurlkurz}{https://devbox42.github.io/sciencesim/spanisch/kl09/unidad-02-mi-nueva-casa/DS-03/LP-03-repaso-general.html}

\begin{document}

\begin{center}

% Titel
\vspace*{10mm}
{\fontsize{28}{34}\selectfont\bfseries\color{spanisch}Lernpfad: \lptitel}

\vspace{15mm}

% QR-Code mit Rahmen
\fbox{\qrcode[height=100mm]{\lpurl}}

\vspace{12mm}

% URL (vollständig, klickbar)
{\small\color{gray}\href{\lpurl}{\texttt{devbox42.github.io/sciencesim/spanisch/kl09/\\unidad-02-mi-nueva-casa/DS-03/LP-03-repaso-general.html}}}

\vspace{8mm}

% Hinweis
{\Large\color{gray}Scanne den QR-Code mit deinem Gerät}

\vspace{10mm}

% Info-Box
\fcolorbox{spanisch}{spanisch!10}{%
    \parbox{0.7\textwidth}{%
        \centering
        \vspace{3mm}
        Klausurvorbereitung U1 + U2 -- Stationenlernen
        \vspace{3mm}
    }%
}

\end{center}

\end{document}
